\documentclass[12pt]{article}
\usepackage[utf8]{inputenc}

%%% I added a few packages. This next one allows you to set the margins.
\usepackage[top=.5in,bottom=1in,right=.75in,left=.75in]{geometry}
\usepackage{multicol}
%\usepackage{newspaper}
\setlength{\columnsep}{1cm}
\usepackage{blindtext}
\usepackage{varwidth}

%%% This one does boxes in a really nice way.
\usepackage{tcolorbox}
\usepackage{mdframed}
\usepackage{minibox}
\usepackage{graphicx}
\graphicspath{ {./images/} }
\usepackage{epstopdf}
  \epstopdfDeclareGraphicsRule{.pdf}{png}{.png}{convert #1 \OutputFile}
  \DeclareGraphicsExtensions{.png,.pdf}
\usepackage{hyperref}
\hypersetup{
    colorlinks=true,
    linkcolor=blue,
    filecolor=magenta,      
    urlcolor=cyan,
    pdftitle={MSCSMess},
    pdfpagemode=FullScreen,
    }
\urlstyle{same}
 \usepackage{fix-cm}
 \usepackage{soul}
 \usepackage{graphicx}
 \usepackage{multicol}


\newcommand{\ds}{\displaystyle}
\newcommand{\ZZ}{\mathbb{Z}}
\newcommand{\QQ}{\mathbb{Q}}
\newcommand{\RR}{\mathbb{R}}
\newcommand{\CC}{\mathbb{C}}
\newcommand{\FF}{\mathbb{F}}
\newcommand{\xx}{\mathbf{x}}
\newcommand{\uu}{\mathbf{u}}
\newcommand{\vv}{\mathbf{v}}

\newcommand{\bolder}{\textbf}
 
 
\title{\fontsize{60}{70}\selectfont {\bf {MSCS} 
\includegraphics[scale=.3]{images/OleLion.png} {Mess}}}
\author{Department of Mathematics, Statistics, and Computer Science \\ St.\@ Olaf College, Northfield, MN 55057 \vspace{-2ex}} 
\date{\today\, $|$ Volume 54, No.\@ 11 \line(1,0){500}}


\begin{document}

\maketitle \vspace{-0.5cm}

Happy Friday! Check out what’s coming up in the MSCS department. As we head into Thanksgiving break, make sure to take time to relax, recharge, and enjoy the holiday. Wishing you a wonderful break filled with rest, good food, and good company!


 \begin{multicols}{2}

 \begin{tcolorbox}[colback=white, colframe=black]
    \begin{center}
    \title{{\large \bolder{MSCS Colloquium 12/1!
}}}
    \end{center}
\end{tcolorbox}

Join Alexander Wiedemann, a Computer Science faculty candidate, who will be giving an MSCS Department Colloquium talk titled “Algorithmics and Complexity of Genome Rearrangement” on Monday, December 1, 2025 at 3:30 PM in RNS 210! Wiedemann’s research sits at the intersection of theory and computation in network and data science, with current work spanning randomized and parameterized algorithms for genome rearrangement, probabilistic inference for opinion dynamics, and public interest data science with community partners. 

In his talk, he will explore the challenges of modeling genome rearrangements, which are large scale mutations that reorder genes within a genome, and discuss how balancing biological realism with computational feasibility leads to deep questions about counting, inference, and complexity. He will also highlight state of the art approaches for several genome rearrangement models, including recent results for the Single Cut and Join model. Everyone is welcome to attend!

 \begin{tcolorbox}[colback=white, colframe=black]
    \begin{center}
    \title{{\large \bolder{Linux Ladies Career Workshop!
}}}
    \end{center}
\end{tcolorbox}

Join the Linux Ladies on Monday, December 2nd from 3–4 PM in RNS 301 for a relaxed, student-focused conversation about career exploration and networking in tech! Come connect with peers, ask questions, and hear insights from fellow students navigating their own career paths! Snacks will be provided!


\begin{tcolorbox}[colback=white, colframe=black]
    \begin{center}
    \title{{\large \bolder{MSCS student profile in the St. Olaf Magazine
}}}
    \end{center}
\end{tcolorbox}

There is an excellent article about Jessica Schmidt--math/bio major, SDS/math bio concentrator, class of 2026--in the current issue of the \href{https://wp.stolaf.edu/magazine/}{St. Olaf Magazine}, titled ``The Cipher Protocol." Jessica is a ``master codebreaker" after having been the only individual to have solved all three codebreaking challenges in the annual Kryptos competition. Read the article (and the follow-up article on the MSCS department and Math program) and congratulate Jessica.

\begin{tcolorbox}[colback=white, colframe=black]
    \begin{center}
    \title{{\large \bolder{CSCI 121 Principles of Computer Science (J-term)
}}}
    \end{center}
\end{tcolorbox}

Dive into computer science with Professor Houssou! This introductory, hands-on course will develop your analytical thinking and problem-solving skills through daily programming exercises. You will apply programming concepts to real-world projects across various fields like science and engineering, as well as music, sociology, finance, art, and more. 

The smaller class size during J-term allows you a more personalized experience into the realm of computer science with one-on-one support from Professor Houssou and close collaboration with your classmates. This is also a great opportunity to build up your technology skills in an ever-increasing technical world.

\begin{tcolorbox}[colback=white, colframe=black]
    \begin{center}
    \title{{\large \bolder{DAAD RISE Internships}}}
    \end{center}
\end{tcolorbox}

Are you interested in a summer STEM internship/fellowship program for undergraduates in Germany, requiring no German-language skills? Check out the DAAD Research Internships in Science and Engineering (RISE) website \href{https://www.daad.de/rise/en/}{at this link}. The application window is October 15th - November 30th.


\begin{tcolorbox}[colback=white, colframe=black]
    \begin{center}
    \title{{\large \bolder{Problem Solving Group}}}
    \end{center}
\end{tcolorbox}
The Math Problem Solving Group meets weekly to have fun, solve math problems, eat candy, and prepare for upcoming math contests. Meetings are 11:30am-12:30pm, in RNS 207 (the Center for Interdisciplinary Research room). Everyone is welcome, regardless of prior experience in mathematics or math contests!


\end{multicols} \vspace{0.5cm}

\vfill

\begin{center}
\includegraphics[width=0.85\textwidth]{Colloquium Alex Wiedemann (2).png}
\end{center}


\begin{center}
\includegraphics[width=0.98\textwidth]{Career Workshop w Linux Ladies.png}
\end{center}


\begin{center}
\includegraphics[width=0.98\textwidth]{images/CSCI 121 J-term Course.png}
\end{center}

\begin{center}
\includegraphics[width=0.9\textwidth]{images/Tea Time Poster.png}
\end{center}

\noindent{To submit an article, event, or anything else for publication in the Mess, email messeditor@stolaf.edu. Visit \href{https://wp.stolaf.edu/mscs/mscs-mess/}{http://wp.stolaf.edu/mscs/mscs-mess/} for a digital archive of previous MSCS Mess issues.}\\
	
\hfill Nima Dahir, Editor
	
\hfill Kyle Salois, Faculty Advisor


 

\end{document}